\subsection{Implementation rules at finite size}\label{sec:rules}
Theorem~\ref{thm:upper} and Proposition~\ref{prop:axis} are asymptotic statements. At finite
size two systematic biases require correction, and both admit analytic expressions without
fitted parameters.

The first is the threshold center shift. The critical threshold shift given by
Dembo--Montanari\cite{dm}, rescaled to the scale of this paper, replaces $z(n)$ with
\begin{equation*}
z_{\rm sh} = z(n) + C_k M^{-1/6}.
\end{equation*}
The constant $C_k$ is given in closed form by Airy integrals and the critical constants
(Appendix~\ref{app:ck}); for $k = 3, 4, 5, 6$, $C_k \approx 1.930$, $2.015$, $2.016$, $2.005$.
The empirical basis is given in Section~\ref{sec:threshold}.

The second is the chi-square shape correction. At finite size the unconditional measurement
is closer to the standardized chi-square $(M-1)$ distribution (Section~\ref{sec:chisq}), and
the quantile requires the additional term
\begin{equation*}
\kappa(u) = \frac{\chi^2_{M-1,u} - (M-1)}{\sqrt{2(M-1)}} - \Phi^{-1}(u).
\end{equation*}
Both corrections vanish as $M \to \infty$ and do not affect the asymptotic coverage.

Both corrections act on the acceptance threshold of Section~\ref{sec:upper-construct}:
replacing $z(n)$ by $z_{\rm sh}$ and adding $\kappa(u)$ to the quantile gives the corrected
selection threshold
\begin{equation}
n + s(n)\bigl[q_{\rm sel}\bigl(u;\,z_{\rm sh},\,\eta_k b^2/n^2\bigr) + \kappa(u)\bigr].
\label{eq:corrected}
\end{equation}
It describes macro-core failure. On the low-load side below the critical window, failure comes
mainly from repeated edges, and using \eqref{eq:corrected} alone gives undercoverage
(Section~\ref{sec:num-loadaxis}). The two implementation rules below handle this in different
ways; the comparison baseline is the unconditional threshold $n + s(n)\Phi^{-1}(u)$, which
applies no selection correction at all.

The envelope rule takes, for each candidate, the lower of \eqref{eq:corrected} and the
unconditional threshold. The acceptance set is the union of the acceptance sets of the two
thresholds, so a candidate is rejected only when both thresholds reject it, and the probability
that the true value is wrongly rejected is at most the probability under either threshold alone.
For $k = 3, 4$ the quantile of the selection law lies above the Gaussian quantile, so the
envelope rule takes the unconditional threshold; for $k \ge 5$ the corrected selection
threshold is lower, and the envelope rule applies the selection correction.

The mixture rule replaces the selection law in \eqref{eq:corrected} with the mixture law $H_n$
of Section~\ref{sec:loadaxis}, again taking the coordinate $z_{\rm sh}$ and adding $\kappa(u)$;
the macro-layer weight is $\Phi(z_{\rm sh})/(\Phi(z_{\rm sh}) + 1 - e^{-\lambda_n})$, where
$1 - e^{-\lambda_n}$ is the Poisson approximation to the repeated-edge probability. It describes
the two failure mechanisms explicitly and covers the entire load axis. The envelope rule is
simpler in form, while the mixture rule gives a narrower upper bound on the low-load side at
$k=3$. End-to-end checks of the two rules are given in Section~\ref{sec:e2e}.
